% The template needs to be compiled using LuaLaTeX or XeLaTeX

\documentclass{article}
\usepackage{jtdh}
% \definecolor{Mycolor2}{HTML}{00F9DE}
\usepackage[utf8]{inputenc} % allow utf-8 input
\usepackage[slovene]{babel} % Change to english if english paper
\usepackage{url}   % simple URL typesetting
\AtBeginDocument{\urlstyle{APACsame}}
\usepackage[hang,flushmargin]{footmisc} % no indentation in footnotes
\usepackage[colorlinks = true, linkcolor = black, urlcolor  = gray, citecolor = black, anchorcolor = black]{hyperref}
\usepackage{booktabs}       % professional-quality tables
\usepackage{amsfonts}       % blackboard math symbols
\usepackage{nicefrac}       % compact symbols for 1/2, etc.
\usepackage{xcolor}
\usepackage{graphicx}
\usepackage{tabularx}
\usepackage[]{apacite}
\usepackage{enumitem}
\graphicspath{{img/}}
% table styles
\aboverulesep=0ex
\belowrulesep=0ex

%\addbibresource{references.bib} %Imports bibliography file 

\title{Naslov prispevka}

\author{
  Ime PRIIMEK,\textsuperscript{1} Ime PRIIMEK\textsuperscript{1, 2}
}
\institutions{
  \textsuperscript{1}Ustanova avtorja \par
  \textsuperscript{2}Ustanova avtorja
}

\begin{document}
\maketitle

\begin{abstract}
Priporočeni obseg izvlečka je od 10 do 15 vrstic. Izvlečku sledi od 3 do 5 ključnih besed v slovenščini (oz. jeziku, v katerem je prispevek napisan).
\end{abstract}
\keywords{ključna beseda 1, ključna beseda 2, ključna beseda 3}


\section{Številčenje naslovov, črkopis, velikost črk, razmik, poravnava, obseg}
Številčenje naslovov in podnaslovov je desetiško (1, 1.1, 1.1.1). Zaželeno je, da členitev nima več kot treh nivojev). Priporočena dolžina prispevka je navedena na spletni strani konference.


\subsection{Navedba avtorja}
Sklic na vir v besedilu je naslednji: \cite{leech1992corpora}. Stran, na kateri se nahaja navedek v delu, se napiše: \cite[str.~107]{leech1992corpora}. Če sta avtorja navedenega dela dva, navedemo oba: \cite{studije2005korpusnem}, pri večjem številu avtorjev izpišemo le prvo ime: \cite{biber1998corpus}, medtem ko bodo v točki Literatura na koncu prispevka načeloma navedeni vsi. Dela enega avtorja, ki so izšla istega leta, med seboj ločimo z zaporednim dodajanjem malih črk (a, b, c itn.) stično ob letnici izida: \cite{erjavec-multext}. Dela različnih avtorjev, ki se nanašajo na isto vsebino, naštejemo po abecednem vrstnem redu, med njimi je podpičje: \cite{erjavec-infotheca, tei}. Če datuma ni (npr. pri spletnih straneh), navedemo na sledeč način: \cite{creativecommons}. 

Nekaj zgledov citiranja:

    \begin{itemize}
        \item Knjiga: \cite{biber1998corpus, studije2005korpusnem}
        \item Članek v znanstveni reviji: \cite{biber1996investigating, erjavec-infotheca}
        \item Članek v monografiji ali konferenčnem zborniku: \cite{leech1992corpora, erjavec-multext}
        \item Slovar: \cite{sinclair1987collins, sskj}
        \item Spletna stran: \cite{openwebspider, creativecommons, scott2008wordsmith, tei}
    \end{itemize}

Primer narativnega citata: \citeA{biber1998corpus} pravijo \ldots 

Če je delo dostopno na spletu, je obvezno navesti njegov URL, ki mora biti stavljen kot hiperlink. Če obstaja, naj se kot povezavo navede trajni identifikator (DOI ali Handle URL). Hiperlinki, ki niso trajni identifikatorji, imajo naveden datum dostopa v obliki: Pridobljeno 23. januarja 2012, \url{http://bos.zrc-sazu.si/sskj.html}. Vsako enoto v teh seznamih zaključuje pika, razen če se konča s hiperlinkom.

\subsection{Citiranje jezikovnih virov} 
Jezikovni viri, kot so korpusi, računalniški leksikoni, jezikovni modeli itd., se navajajo v razdelku za literaturo, in sicer na enak način kot knjige, kar pomeni, da so navedeni tudi avtorji. Repozitorij, kjer je korpus deponiran, z vidika citiranja ustreza založniku pri knjigah; glej \cite{parliamentaryCorpus}. 

Če se isti vir nahaja na več lokacijah, navedemo različico, ki smo jo dejansko uporabili. Na primer: \shortcite{gigafida-clarin} v primeru korpusa \textit{Gigafida 2.0}, če ste ga uporabili preko konkordančnika CLARIN.SI, oziroma \shortcite{gigafida-cjvt}, če ste ga uporabili preko konkordančnika CJVT.

Če za vir obstaja trajni identifikator, navedemo le-tega: \url{http://hdl.handle.net/11356/1748} pri korpusu \textit{siParl 3.0} \cite{parliamentaryCorpus}.\footnote{Velja opozoriti, da \url{https://www.clarin.si/repository/xmlui/handle/11356/1748} ni trajni identifikator.}





\subsection{Opombe, izpusti in daljši navedki}
Številko za opombo vnesemo za ločilom.\footnote{To je primer za besedilo opombe.} Če je v opombi samo URL-naslov, na koncu ni pike.

Izpust iz navedka označimo s poševnicami in tremi pikami: /\ldots/. Na začetku in na koncu navedka oznaka izpusta ni potrebna. Daljši navedki (več kot 3 vrstice) naj bodo postavljeni v samostojen odstavek z odmikom od levega roba. Taki navedki so brez narekovajev. Pred umaknjenim odstavkom ni prazne vrstice, za njim prav tako ne. Primer daljšega navedka:

\begin{quote}
Da se je Ramovš zavzel za vpeljavo teh oblik – četudi pogojno – v oficielni pravopis, pa je bila posledica njegovega novega, samostojnega študija slovenskega jezika, pri katerem je odkril takoimenovani kratki nedoločnik kot prastaro, še predtrubarsko varianto dolgega nedoločnika \cite[str.~198]{vodusek1959}.
\end{quote}

\subsection{Navedba tabel in slik}
\subsubsection{Številčenje in naslavljanje tabel ter slik}
Tabele in slike številčimo zaporedno. Nanje se sklicujemo z rabo velike začetnice (npr. v Tabeli \ref{slogi} prikazujemo …). Naslov je nad tabelo ali sliko, zaključuje ga pika. Besedilne vrednosti so v tabeli poravnane levo, številčne pa desno. Za morebitno ločevanje tisočic rabimo piko, za decimalke pa vejico

\begin{table}[h]
\caption{Naslov tabele.}
\label{slogi}
\begin{tabularx}{\textwidth}{X|X|r}
\toprule
\textit{Kategorija 1} & \textit{Kategorija 2} & \textit{Kategorija 3 –   številčne vrednosti} \\ \midrule
Prva vrstica          & Besedilo tabele       & 1.900,12                                      \\
Druga vrstica         & Besedilo tabele       & 13.934,34                                     \\ \bottomrule
\end{tabularx}
\end{table}

\subsection{Posebni slogi}
Priporočeni slog za alineje je naslednji:
\begin{itemize}
    \item seznam, postavka 1;
    \item seznam, postavka 2.
\end{itemize}

\acknowledgment{
Primer besedila v zahvali.
}

\bibliographystyle{apacite}
\urlstyle{same}
\bibliography{references}

% \bibliography{references}  %%% Remove comment to use the external .bib file (using bibtex).
%%% and comment out the ``thebibliography'' section.
%\printbibliography

\newpage
\titleeng{Naslov članka v angleščini}
\begin{abstracteng}
Na lastni strani na koncu prispevka je naslov članka in povzetek v angleškem jeziku (oz. naslov in povzetek v slovenskem jeziku, če je članek napisan v angleščini). Povzetku sledi tri do pet ključnih besed. Na isti strani (na dnu) je obvestilo o licenci.
\end{abstracteng}

\keywordseng{keyword 1, keyword 2, keyword 3}

\end{document}